\documentclass[12pt]{article}
\usepackage{amsmath, amssymb}
\usepackage{geometry}
\geometry{margin=1in}

\title{LECTURE 19 — N-VARIATE JOINT DISTRIBUTIONS\\
\large High-Density Logical Revision Map}
\author{}
\date{}

\begin{document}
\maketitle

\section{WHY JOINT DISTRIBUTIONS? (Concept Foundation)}

\subsection*{Step 1: Representation}

\[
\mathbf{X} = [X_1, X_2, \dots, X_N]^T
\]

\textbf{Reasoning:}
\begin{itemize}
    \item Instead of treating variables separately, we group them.
    \item Joint distribution describes probability behavior of all variables together.
    \item Transition: single-variable $\rightarrow$ multi-variable framework.
\end{itemize}

\subsection*{Step 2: Forms of Representation}

\begin{itemize}
    \item Joint PMF (discrete)
    \item Joint PDF (continuous)
    \item Joint CDF
\end{itemize}

\textbf{Reasoning:}
\begin{itemize}
    \item Direct extensions of single-variable definitions.
    \item Now we consider simultaneous events.
\end{itemize}

\section{CORE DEFINITIONS}

\subsection{Joint PMF}

\[
P_{X_1,\dots,X_N}(x_1,\dots,x_N) = P(X_1 = x_1, \dots, X_N = x_N)
\]

\textbf{Reasoning:}
\begin{itemize}
    \item Probability of all variables taking values simultaneously.
    \item Represents one combined event.
\end{itemize}

\subsection{Joint CDF}

\[
F(x_1,\dots,x_N) = P(X_1 \le x_1, \dots, X_N \le x_N)
\]

\textbf{Reasoning:}
\begin{itemize}
    \item Accumulates probability over an N-dimensional region.
    \item Extension of 1D CDF: line $\rightarrow$ region.
\end{itemize}

\subsection{Joint PDF}

\[
f = \frac{\partial^N F}{\partial x_1 \cdots \partial x_N}
\]

\textbf{Reasoning:}
\begin{itemize}
    \item Same logic as 1D: PDF = derivative of CDF.
    \item For $N$ variables, take $N$ partial derivatives.
\end{itemize}

\section{INDEPENDENCE CONDITION}

\[
f(x_1,\dots,x_N) = \prod_{i=1}^{N} f_{X_i}(x_i)
\]

\textbf{Reasoning:}
\begin{itemize}
    \item Joint PDF equals product of marginals.
    \item This is the test for independence.
\end{itemize}

\section{EXAMPLE 1 — JOINT GAUSSIAN}

\subsection*{Step 1: Given}

\[
\mu_X = 2, \quad \sigma_X^2 = 1
\]
\[
\mu_Y = -3, \quad \sigma_Y^2 = 4
\]
\[
\text{cov}(X,Y) = -1
\]

\subsection*{Step 2: Correlation}

\[
\rho = \frac{\text{cov}(X,Y)}{\sigma_X \sigma_Y}
\]

\textbf{Reasoning:} Gaussian formula uses correlation.

\subsection*{Step 3: Standard Deviations}

\[
\sigma_X = 1, \quad \sigma_Y = 2
\]

\subsection*{Step 4: Compute $\rho$}

\[
\rho = \frac{-1}{(1)(2)} = -\frac{1}{2}
\]

\subsection*{Step 5: Joint PDF}

Substitute values into Gaussian formula.

\subsection*{Step 6: Marginals}

\[
f_X(x) = \int f_{XY}(x,y)\,dy, \quad
f_Y(y) = \int f_{XY}(x,y)\,dx
\]

\[
X \sim N(\mu_X, \sigma_X^2), \quad Y \sim N(\mu_Y, \sigma_Y^2)
\]

\textbf{Reasoning:} Marginals of jointly Gaussian variables are Gaussian.

\subsection*{Step 7: Compute $P(X < 0)$}

\[
P(X < 0) = F_X(0)
\]

\subsubsection*{Standardization}

\[
Z = \frac{X - \mu_X}{\sigma_X}
\]

\[
P(X < 0) = P\left(Z < \frac{0 - 2}{1}\right) = P(Z < -2)
\]

\subsubsection*{Symmetry}

\[
\Phi(-x) = 1 - \Phi(x) = P(Z > x)
\]

\[
P(Z < -2) = P(Z > 2)
\]

\subsubsection*{Q-function}

\[
Q(x) = P(Z > x)
\]

\[
P(X < 0) = Q(2)
\]

\section{APPLICATION — LARGE MIMO SYSTEM}

\subsection*{System}

\begin{itemize}
    \item $M$ transmit antennas
    \item $N$ receive antennas
\end{itemize}

\subsection*{Channel Model}

\[
h_{ij} = \text{random variable}
\]

\subsection*{Matrix Form}

\[
H =
\begin{bmatrix}
h_{11} & \cdots & h_{1N} \\
\vdots & \ddots & \vdots \\
h_{M1} & \cdots & h_{MN}
\end{bmatrix}
\]

\textbf{Reasoning:} Multiple random variables $\Rightarrow$ joint modeling.

\section{EXAMPLE 2 — UNIFORM 3D VECTOR}

\subsection*{Given}

\[
|x_1| \le 1, \quad |x_2| \le 1, \quad |x_3| \le 1
\]

\[
f =
\begin{cases}
C, & \text{inside cube} \\
0, & \text{otherwise}
\end{cases}
\]

\subsection*{Find $C$}

\[
\iiint f = 1
\]

\[
C (2)(2)(2) = 1 \Rightarrow C = \frac{1}{8}
\]

\subsection*{Marginals}

\[
f_{X_1,X_2} = \int_{-1}^{1} \frac{1}{8} dx_3 = \frac{1}{4}
\]

\[
f_{X_1} = \int_{-1}^{1} \int_{-1}^{1} \frac{1}{8} dx_2 dx_3 = \frac{1}{2}
\]

\subsection*{Conditional PDFs}

\[
f_{X_1|X_2,X_3} = \frac{1/8}{1/4} = \frac{1}{2}
\]

\[
f_{X_1,X_2|X_3} = \frac{1/8}{1/2} = \frac{1}{4}
\]

\subsection*{Independence Check}

\[
\text{LHS} = \frac{1}{8}
\]

\[
\text{RHS} = \frac{1}{2} \cdot \frac{1}{2} \cdot \frac{1}{2} = \frac{1}{8}
\]

\textbf{Conclusion:} Independent

\section{EXAMPLE 3 — TRIVARIATE EXPONENTIAL}

\subsection*{Given}

\[
f_{XYZ} = K e^{-(ax + by + cz)}, \quad x,y,z > 0
\]

\subsection*{Find $K$}

\[
\iiint f = 1
\]

\[
e^{-(ax+by+cz)} = e^{-ax} e^{-by} e^{-cz}
\]

\[
K \left(\frac{1}{a}\right)\left(\frac{1}{b}\right)\left(\frac{1}{c}\right) = 1
\]

\[
K = abc
\]

\subsection*{Marginals}

\[
f_{XY} = ab e^{-(ax+by)}
\]

\[
f_X = a e^{-ax}
\]

\subsection*{Independence Check}

\[
\text{LHS} = abc e^{-(ax+by+cz)}
\]

\[
\text{RHS} = (ae^{-ax})(be^{-by})(ce^{-cz}) = abc e^{-(ax+by+cz)}
\]

\textbf{Conclusion:} Independent

\section{FINAL REVISION FLOW (MENTAL MAP)}

\begin{itemize}
    \item \textbf{Vector Form} $\rightarrow$ Joint Behavior
    
    \item \textbf{PMF / CDF / PDF} $\rightarrow$ Core Definitions
    
    \item \textbf{Independence} $\rightarrow$ Product Rule
    
    \item \textbf{Gaussian Case:}
    \begin{itemize}
        \item $\text{cov} \rightarrow \rho$
        \item Substitute into joint PDF
        \item Marginals
        \item Standardization
        \item $Q$-function
    \end{itemize}
    
    \item \textbf{Uniform Cube:}
    \begin{itemize}
        \item Normalize
        \item Integrate (marginals)
        \item Conditional PDFs
        \item Independence check
    \end{itemize}
    
    \item \textbf{Exponential Case:}
    \begin{itemize}
        \item Factorization
        \item Normalize
        \item Marginals
        \item Independence check
    \end{itemize}
    
\end{itemize}

\end{document}
